\chapter{Results}\label{ch:results}

\section{Exploratory Summaries}
Temperature and power have a positive Pearson correlation\index{correlation coefficient} in this generated dataset ($r=\rTP$; \cref{fig:res-panels}). The hinge regression has an in-sample coefficient of determination\index{coefficient of determination} of $R^2=\hingeRtwo$. This is a summary of the same rows used to fit the model, not an estimate of performance on unseen measurements.

\Cref{fig:res-panels} contains four views of the synthetic data: the daily power series in panel (a), the variable-correlation matrix in panel (b), the site distributions in panel (c), and power against temperature in panel (d). The final panel shows both the fitted hinge curve and a dashed linear baseline. Per-site counts and summary values are listed in \cref{tab:res-summary}.

\begin{figure}[htbp]
  \centering
  \begin{subfigure}[t]{0.48\textwidth}
    \includegraphics[width=\linewidth]{fig_panel_a.pdf}
    \caption{Daily mean power.}
  \end{subfigure}\hfill
  \begin{subfigure}[t]{0.48\textwidth}
    \includegraphics[width=\linewidth]{fig_panel_b.pdf}
    \caption{Correlation matrix.}
  \end{subfigure}\\[4pt]
  \begin{subfigure}[t]{0.48\textwidth}
    \includegraphics[width=\linewidth]{fig_panel_c.pdf}
    \caption{Power by site.}
  \end{subfigure}\hfill
  \begin{subfigure}[t]{0.48\textwidth}
    \includegraphics[width=\linewidth]{fig_panel_d.pdf}
    \caption{Power against temperature.}
  \end{subfigure}
  \caption{Four summaries generated from the synthetic hourly dataset. Panel (d) shows the fitted hinge curve and the linear baseline.}
  \label{fig:res-panels}
\end{figure}

\begin{table}[htbp]
  \centering
  \caption{Summary statistics by example site.}
  \label{tab:res-summary}
  \input{../../figures/summary_table}
\end{table}

\section{Discussion}
Site B has the highest mean power in this generated sample. The data-generation script adds a positive site offset\index{site offset} to Site B, so this difference is built into the example. The correlation and $R^2$ describe the generated rows; they do not establish a causal relationship or validate a forecast for a real site.

Proposition \cref{prop:ols} concerns the unbiasedness and covariance of the ordinary least-squares estimator under its stated assumptions. It is a result about the estimator, not evidence that the hinge specification is the correct model for real measurements.
